% ---------- Слайд 2: Что такое LaTeX ----------
\begin{frame}[label=what-is-latex]{Что такое \LaTeX}
  \begin{block}{Определение}
    \LaTeX  --- система компьютерной вёрстки, ориентированная
    на подготовку научных документов.
  \end{block}

  \begin{itemize}
    \item Работает по принципу "исходный код $\to$ скомпилированный документ";
    \item Единый инструмент для статей, диссертаций и слайдов;
    \item Основан на \TeX{} (Д.~Кнут, 1978), развивается с 1984~г.
  \end{itemize}
  \note{Историческую справку не раскрывать. Основной тезис --- единый инструмент для всех типов научных документов.}
\end{frame}

% ---------- Слайд 3: WYSIWYG vs WYSIWYM ----------
\begin{frame}[label=wysiwyg]{WYSIWYG и WYSIWYM}
  \begin{columns}[T,onlytextwidth]
    \column{0.48\textwidth}
      \begin{block}{WYSIWYG (Word, PowerPoint)}
        What You See Is What You Get
        \begin{itemize}
          \item Результат виден в процессе набора;
          \item Оформление задаётся вручную;
          \item Стиль изменяется при редактировании;
          \item Формулы требуют отдельных решений.
        \end{itemize}
      \end{block}
    \column{0.48\textwidth}
      \begin{exampleblock}{WYSIWYM (\LaTeX)}
        What You See Is What You Mean
        \begin{itemize}
          \item Описывается логическая структура документа;
          \item Оформление задаётся автоматически;
          \item Стиль сохраняется при любых правках;
          \item Формулы --- штатный элемент разметки.
        \end{itemize}
      \end{exampleblock}
  \end{columns}

  \vspace{0.4em}
  \centering
  \emph{В \LaTeX{} содержание отделено от оформления.}
  \note{Ключевое различие: в \LaTeX{} пользователь описывает структуру,
  а типографика обеспечивается системой.}
\end{frame}

% ---------- Слайд 4: От статьи к слайдам ----------
\begin{frame}[label=article-to-slides]{Как статьи, так и слайды}
  \begin{block}{Универсальность}
    Наличие опыта написания статей в \LaTeX{} практически гарантирует простоту освоения презентаций.
  \end{block}

  \begin{itemize}
    \item Формулы, таблицы, \cmd{cite}, \cmd{includegraphics} работают идентично;
    \item Различие состоит в классе документа: \texttt{article} $\to$ \texttt{beamer};
    \item Переход к презентациям требует освоения лишь небольшого набора команд.
  \end{itemize}
  \note{Слайд важен для аудитории, уже использующей \LaTeX{} в статьях: оценить объём дополнительных знаний.}
\end{frame}